% Discussion
% Three subsections: failure analysis, key insights and practical implications,
% future work.

\section{Discussion}
\label{sec:discussion}

\subsection{Failure Analysis}
\label{sec:disc-failures}

In this section we conduct a manual qualitative analysis to better understand the misclassification patterns across the explored IRC methods. We inspect the misclassifications produced by Qwen-32B (the strongest model) under each method's best configuration. From all the misclassifications across \ragtag. \bragtag, and fine-tune, we drew a random sample of 50 misclassifications per method (150 total). Each Method's sample includes 30 true-question, 10 true-bug, and 10 true-feature errors. Each case was analyzed by two authors independently, with discrepancies resolved through discussions. The observed failure categories and their frequencies are as follows:

\textbf{\textit{Wrong template format (77.33\%).}} Issue reports whose content and ground-truth label contradicts their chosen template and structure. All of these errors are from questions and feature requests that the user filed under a bug-report template (\textit{Steps to reproduce}, \textit{Expected}, \textit{Actual}, \textit{System Information}). Consequently, the structure of the bug template misleads the LLM classifier towards predicting bug, regardless of the issue's content. For example, \texttt{microsoft/vscode\#193985}\footnote{\url{https://github.com/microsoft/vscode/issues/193985}} is an actual support question that opens with a literal \texttt{Type:Bug} because it uses VS Code's bug template.

\textbf{\textit{Hybrid intent (13.33\%).}} Issues that could have multiple defensible labels due to their combined intents. These errors mostly take place when a fix or a feature request is layered in a question or bug. For example, in \texttt{kubernetes/kubernetes\#120364}\footnote{\url{https://github.com/kubernetes/kubernetes/issues/120364}} the author both reports a concrete bug and in the same body explicitly requests a new feature: ``\textit{Please, give us multiarchitecture image! and fix monoarchitecture images too!}''

\textbf{\textit{Label noise (9.33\%).}} Issues whose ground-truth label does not match its content. For example, \texttt{tensorflow/tensorflow\#61710}\footnote{\url{https://github.com/tensorflow/tensorflow/issues/61710}} is titled ``\textit{Will there a MLP model in the future version?}'' which requests a new API feature for building MLPs. Yet the dataset labels it as a bug.


We also sampled 30 issues with \textit{invalid} predictions per method (60 total) and observed three failure modes. In roughly 73\% of cases, the model either continued to generate issue-body content or its own reasoning. Another 23\% were hallucinations such as repeated digits and punctuations. The remaining 3\% are parsable labels that fall outside the three valid classes (e.g., \texttt{documentation}).


\subsection{Key Insights and Practical Implications}
\label{sec:disc-insights}

\textbf{The practicality of context retrieval for IRC.} Our results indicate that RAG-based approaches offer a more practical solution for IRC than fine-tuning. With $11\times$ less labeled data per project and 25--47\% less peak GPU memory, \ragtag\ performs competitively against fine-tuning (trails by 0.020 in aggregate macro F1), and \bragtag\ closes the remaining gap to statistical equivalence, passing TOST equivalence at $\delta{=}0.01$ (\Cref{sec:method-comparison}). RAG-based methods offer deployment advantages over fine-tuning as well: adding a new labeled issue report becomes an incremental index update rather than a retraining cycle, and the model can be swapped in production with no project-specific retraining. For extremely resource-constrained scenarios, even \votag PS is a defensible baseline: it reaches macro $F_1$ = 0.595 ($k{=}15$), within 0.018 of Qwen-3B's zero-shot 0.613 (\Cref{sec:rq1}, \Cref{sec:ragtag}), at a fraction of the GPU cost ($\approx$240~MB vs 2.9~GB).

\textbf{Combining classification and validation.} Our failure analysis (\Cref{sec:disc-failures}) shows that users frequently file support questions and feature requests under the bug-report template, which can bias the classifier towards bug regardless of the content. This emphasizes the importance of refining the new LLM-assisted triage pipelines now in place on platforms such as GitHub~\cite{githubcopilot} and GitLab~\cite{gitlabduo}. In these pipelines, the reporter describes the report to an LLM, which then suggest the correct template and labels. Report validation is also important since mislabeling and duplication is a common problem among contributors\cite{antoniol2008bug, pandey2017automated}. Studies have already explored validation for bug reports~\cite{khatib2025assertflip, wang2024aegis, ahmed2025otter, dincc2025judge}. However, combining classification and validation in a cohesive workflow is underexplored.

\subsection{Future Work}
\label{sec:disc-future}

\textbf{Generalization to more diverse environments.} Our analysis covers a large benchmark across 11 GitHub projects with a balanced 600 issues per project. We encourage future work to evaluate IRC methods under more diverse settings such as imbalanced label distributions with one or more minority classes, multi-label settings where more than one label can be predicted for a single issue, and other issue tracking platforms (e.g., Jira, Bugzilla). These future directions will assess the generalizability of our findings and uncover new patterns.

\textbf{Other retrieval techniques.} Our study empirically shows that retrieval heuristics can soften model biases and even match fine-tuning without parameter updates. This opens the door for future research on other retrieval-side interventions for different IRC settings. A concrete future direction is a two-stage inference scheme where a first LLM call filters or re-ranks the retrieved neighbors based on their relevance, and a second call performs the classification. Settings such as embedding model and similarity metric are also worth exploring in future work, because of their impact on retrieval quality.

